% GENERATED FILE. Edit canonical lessons, then run npm run generate:books.
% canonical-source-hash: fnv1a64:ca9aa6908f9b95fe
% canonical-lessons: RU-C526-oblozhka, RU-C526-zakladka, RU-C526-listat, RU-C526-izdatelstvo, RU-C526-perechityvat

\chapter{Covers and Bookmarks}
\label{ch:ru-covers-and-bookmarks}
\hlchaptermodality{526}

\begin{chapteropening}
\textbf{By the end of this chapter:} \emph{I can talk about handling and rereading books.}

Complete the last lesson of chapter 526: I can talk about handling and rereading books.
\end{chapteropening}

\section[\texorpdfstring{\ru{обложка} (oblózhka)}{обложка (oblózhka)}]{\ru{обложка} (oblózhka) — a cover (of a book)}

\label{lesson:RU-C526-oblozhka}



\begin{quote}
Before the new one: say the Russian for to rhyme, then the Russian for by heart.
\end{quote}

\subsection*{You'll want to know: \ru{обложка}}
\textbf{\ru{обложка}} (\emph{oblózhka}) — \textquotedblleft{}a cover (of a book)\textquotedblright{}.

\begin{quote}
\ru{Мне} \ru{понравилась} \ru{обложка} \ru{этой} \ru{книги}.
\end{quote}

\begin{quote}
\emph{Mne ponrávilas oblózhka étoy knígi.}
\end{quote}

\begin{quote}
I liked the cover of this book.
\end{quote}

Feminine.

\begin{cousinweb}
\emph{\ru{обложка}} is built on \emph{\ru{обложить}} 'to cover, to lay round', from the root of \emph{\ru{ложить}} 'to lay'; EtymDB has no entry.
\end{cousinweb}

\begin{grammarlens}[title={What you've built}]
One more word for telling, explaining and describing at B1.
\end{grammarlens}

\subsection*{Your turn}
\begin{itemize}
\raggedright
  \item \emph{Say it:} \emph{oblózhka}
  \item \emph{Say it:} \emph{oblózhka}, once more
  \item \emph{Say it:} a sentence of your own with \emph{oblózhka}
  \item \emph{Recall:} say the Russian for to rhyme, then the Russian for by heart, then say \emph{oblózhka} again
\end{itemize}

\begin{tcolorbox}[breakable,colback=teal!4,colframe=teal!35!black,title={Before you move on}]
What does \ru{обложка} mean? (\textquotedblleft{}A cover (of a book)\textquotedblright{}.) Say it once more.
\end{tcolorbox}
\section[\texorpdfstring{\ru{закладка} (zakládka)}{закладка (zakládka)}]{\ru{закладка} (zakládka) — a bookmark}

\label{lesson:RU-C526-zakladka}



\begin{quote}
Before the new one: say the Russian for by heart, then the Russian for a cover.
\end{quote}

\subsection*{You'll want to know: \ru{закладка}}
\textbf{\ru{закладка}} (\emph{zakládka}) — \textquotedblleft{}a bookmark\textquotedblright{}.

\begin{quote}
\ru{Я} \ru{потерял} \ru{закладку} \ru{и} \ru{не} \ru{помню} \ru{страницу}.
\end{quote}

\begin{quote}
\emph{Ya poteryál zakládku i ne pómnyu stranítsu.}
\end{quote}

\begin{quote}
I've lost my bookmark and can't remember the page.
\end{quote}

Feminine. Also a browser bookmark.

\begin{cousinweb}
\emph{\ru{закладка}} is built on \emph{\ru{закладывать}} 'to put in, to insert', from \emph{\ru{класть}} 'to put'; EtymDB has no entry.
\end{cousinweb}

\begin{grammarlens}[title={What you've built}]
One more word for telling, explaining and describing at B1.
\end{grammarlens}

\subsection*{Your turn}
\begin{itemize}
\raggedright
  \item \emph{Say it:} \emph{zakládka}
  \item \emph{Say it:} \emph{zakládka}, once more
  \item \emph{Say it:} a sentence of your own with \emph{zakládka}
  \item \emph{Recall:} say the Russian for by heart, then the Russian for a cover, then say \emph{zakládka} again
\end{itemize}

\begin{tcolorbox}[breakable,colback=teal!4,colframe=teal!35!black,title={Before you move on}]
What does \ru{закладка} mean? (\textquotedblleft{}A bookmark\textquotedblright{}.) Say it once more.
\end{tcolorbox}
\section[\texorpdfstring{\ru{листать} (listát)}{листать (listát)}]{\ru{листать} (listát) — to leaf through}

\label{lesson:RU-C526-listat}



\begin{quote}
Before the new one: say the Russian for a cover, then the Russian for a bookmark.
\end{quote}

\subsection*{You'll want to know: \ru{листать}}
\textbf{\ru{листать}} (\emph{listát}) — \textquotedblleft{}to leaf through\textquotedblright{}.

\begin{quote}
\ru{В} \ru{очереди} \ru{я} \ru{листал} \ru{журнал}.
\end{quote}

\begin{quote}
\emph{V ócheredi ya listál zhurnál.}
\end{quote}

\begin{quote}
While I was in the queue I leafed through a magazine.
\end{quote}

Imperfective; perfective \ru{пролистать}.

\begin{cousinweb}
\emph{\ru{листать}} is built on \emph{\ru{лист}} 'a leaf, a sheet of paper'; EtymDB lists the verb without a source.
\end{cousinweb}

\begin{grammarlens}[title={What you've built}]
One more word for telling, explaining and describing at B1.
\end{grammarlens}

\subsection*{Your turn}
\begin{itemize}
\raggedright
  \item \emph{Say it:} \emph{listát}
  \item \emph{Say it:} \emph{listát}, once more
  \item \emph{Say it:} a sentence of your own with \emph{listát}
  \item \emph{Recall:} say the Russian for a cover, then the Russian for a bookmark, then say \emph{listát} again
\end{itemize}

\begin{tcolorbox}[breakable,colback=teal!4,colframe=teal!35!black,title={Before you move on}]
What does \ru{листать} mean? (\textquotedblleft{}To leaf through\textquotedblright{}.) Say it once more.
\end{tcolorbox}
\section[\texorpdfstring{\ru{издательство} (izdátelstvo)}{издательство (izdátelstvo)}]{\ru{издательство} (izdátelstvo) — a publishing house}

\label{lesson:RU-C526-izdatelstvo}



\begin{quote}
Before the new one: say the Russian for a bookmark, then the Russian for to leaf through.
\end{quote}

\subsection*{You'll want to know: \ru{издательство}}
\textbf{\ru{издательство}} (\emph{izdátelstvo}) — \textquotedblleft{}a publishing house\textquotedblright{}.

\begin{quote}
\ru{Его} \ru{роман} \ru{вышел} \ru{в} \ru{большом} \ru{издательстве}.
\end{quote}

\begin{quote}
\emph{Yegó román výshel v bolshóm izdátelstve.}
\end{quote}

\begin{quote}
His novel came out with a big publishing house.
\end{quote}

Neuter.

\begin{cousinweb}
\emph{\ru{издательство}} is built on \emph{\ru{издатель}} 'a publisher', from \emph{\ru{издать}} 'to publish', which is built on \emph{\ru{из}-} and \emph{\ru{дать}} 'to give'; EtymDB lists it, glossed 'publishing house', without a source.
\end{cousinweb}

\begin{grammarlens}[title={What you've built}]
One more word for telling, explaining and describing at B1.
\end{grammarlens}

\subsection*{Your turn}
\begin{itemize}
\raggedright
  \item \emph{Say it:} \emph{izdátelstvo}
  \item \emph{Say it:} \emph{izdátelstvo}, once more
  \item \emph{Say it:} a sentence of your own with \emph{izdátelstvo}
  \item \emph{Recall:} say the Russian for a bookmark, then the Russian for to leaf through, then say \emph{izdátelstvo} again
\end{itemize}

\begin{tcolorbox}[breakable,colback=teal!4,colframe=teal!35!black,title={Before you move on}]
What does \ru{издательство} mean? (\textquotedblleft{}A publishing house\textquotedblright{}.) Say it once more.
\end{tcolorbox}
\section[\texorpdfstring{\ru{перечитывать} (perechítyvat)}{перечитывать (perechítyvat)}]{\ru{перечитывать} (perechítyvat) — to reread}

\label{lesson:RU-C526-perechityvat}



\begin{quote}
Before the new one: say the Russian for to leaf through, then the Russian for a publishing house.
\end{quote}

\subsection*{You'll want to know: \ru{перечитывать}}
\textbf{\ru{перечитывать}} (\emph{perechítyvat}) — \textquotedblleft{}to reread\textquotedblright{}.

\begin{quote}
\ru{Я} \ru{часто} \ru{перечитываю} \ru{любимые} \ru{книги}.
\end{quote}

\begin{quote}
\emph{Ya chásto perechítyvayu lyubímyye knígi.}
\end{quote}

\begin{quote}
I often reread my favourite books.
\end{quote}

Imperfective; perfective \ru{перечитать}.

\begin{cousinweb}
\emph{\ru{перечитывать}} is built on \emph{\ru{пере}-} 'again' and \emph{\ru{читать}} 'to read'; EtymDB has no entry.
\end{cousinweb}

\begin{grammarlens}[title={What you've built}]
One more word for telling, explaining and describing at B1.
\end{grammarlens}

\subsection*{Your turn}
\begin{itemize}
\raggedright
  \item \emph{Say it:} \emph{perechítyvat}
  \item \emph{Say it:} \emph{perechítyvat}, once more
  \item \emph{Say it:} a sentence of your own with \emph{perechítyvat}
  \item \emph{Recall:} say the Russian for to leaf through, then the Russian for a publishing house, then say \emph{perechítyvat} again
\end{itemize}

\begin{tcolorbox}[breakable,colback=teal!4,colframe=teal!35!black,title={Before you move on}]
What does \ru{перечитывать} mean? (\textquotedblleft{}To reread\textquotedblright{}.) Say it once more.
\end{tcolorbox}
